\chapter{Criterio de selección prioritaria de claims}

\section{Selección prioritaria}
\begin{definition}
Se llama selección prioritaria de claims al procedimiento curatorial mediante el cual, dentro de la tabla de claims, ciertas afirmaciones reciben precedencia metodológica para revisión, desarrollo o eventual exposición pública.
\end{definition}

\begin{observation}
La prioridad no equivale todavía a validación final ni a superioridad absoluta. Su función es ordenar el trabajo metodológico y evitar una curación plana de todas las afirmaciones.
\end{observation}

\section{Condición mínima de prioridad}
\begin{definition}
Un claim puede recibir prioridad cuando concurren de manera suficiente una o más de las siguientes condiciones:
\begin{enumerate}
  \item claridad estructural especialmente alta;
  \item trazabilidad particularmente limpia;
  \item relación fuerte entre evidencia disponible y formulación disciplinada;
  \item valor ejemplificador para el corpus;
  \item utilidad metodológica para ordenar otros claims o demos relacionadas.
\end{enumerate}
\end{definition}

\begin{observation}
La prioridad no debe asignarse por atractivo retórico. Debe responder a fuerza metodológica discernible.
\end{observation}

\section{Prioridad y tabla de claims}
\begin{definition}
La tabla de claims debe permitir marcar prioridad editorial o metodológica sin borrar el estado, el alcance o la fragilidad relativa de cada afirmación.
\end{definition}

\begin{observation}
Esta cláusula evita que la prioridad funcione como disfraz de cierre probatorio prematuro.
\end{observation}

\section{Prioridad y auditoría}
\begin{definition}
La selección prioritaria debe ser compatible con la auditoría mínima de claims y no puede ignorar observaciones críticas ya registradas.
\end{definition}

\begin{observation}
Un claim con observación crítica puede seguir siendo prioritario, pero solo si su prioridad queda expresamente motivada y no oculta su fragilidad.
\end{observation}

\section{Condición mínima de no prioridad}
\begin{definition}
Un claim no debe recibir prioridad cuando su formulación sigue siendo demasiado opaca, su evidencia es todavía insuficiente para el propósito declarado o su posición depende de extrapolaciones no disciplinadas.
\end{definition}

\begin{observation}
La no prioridad no equivale a descarte definitivo. Puede significar solo que el claim debe seguir en curación antes de ocupar un lugar central.
\end{observation}

\section{Cláusula de alcance}
\begin{observation}
El criterio de selección prioritaria fija una primera jerarquía curatorial dentro de la tabla de claims. No constituye todavía una selección final pública ni una clasificación exhaustiva del valor de todas las afirmaciones del corpus.
\end{observation}

